\makeatletter
\ifx \nauthor\undefined
\def\nauthor{Anonymous}
\else
\fi

\author{Based on lectures by \nlecturer \\\small Notes taken by \nauthor}
\date{\nterm\ \nyear}

\usepackage{alltt}
\usepackage{mdframed}
\usepackage{amsfonts}
\usepackage{amsmath}
\usepackage{amssymb}
\usepackage{amsthm}
\usepackage{booktabs}
\usepackage{caption}
\usepackage{enumitem}
\usepackage{fancyhdr}
\usepackage{graphicx}
\usepackage{mathdots}
\usepackage{mathtools}
\usepackage{microtype}

\usepackage{listings}

\usepackage{multirow}
\usepackage{pdflscape}
\usepackage{pgfplots}
\usepackage{siunitx}
\usepackage{slashed}
\usepackage{tabularx}
\usepackage{tkz-euclide}
\usepackage[normalem]{ulem}
\usepackage[all]{xy}
\usepackage{imakeidx}
\usepackage{float}
\usepackage{avant}
\usepackage{eso-pic}
\usepackage{tikz}
\usepackage{xcolor}
\usepackage[a4paper,margin=2.5cm]{geometry}
\usepackage{titlesec}
\usepackage[most]{tcolorbox}
\usepackage{tocloft}

% ============================================================
% Proof and Solution environments
% ============================================================

\renewenvironment{proof}[1][]%
{%
    \par
    \pushQED{\qed}%
    \normalfont
    \topsep6\p@\@plus6\p@\relax
    \trivlist
    \item[\hskip\labelsep
        \itshape
        \ifx\relax#1\relax
            Proof%
        \else
            Proof of `#1'%
        \fi
        \@addpunct{.}]%
}%
{\popQED\endtrivlist\@endpefalse}

\newenvironment{solu}[1][]%
{%
    \par
    \normalfont
    \topsep6\p@\@plus6\p@\relax
    \trivlist
    \item[\hskip\labelsep
        \itshape
        \ifx\relax#1\relax
            Solution%
        \else
            Solution to `#1'%
        \fi
        \@addpunct{.}]%
}%
{\endtrivlist\@endpefalse}
%Reformatting Sections and Subsections with these commands

\titleformat{\section}
  {\Large\bfseries\color{durhampurple}}
  {\thesection}
  {0.75em}
  {}

\titleformat{\subsection}
  {\large\bfseries\color{durhamviolet}}
  {\thesubsection}
  {0.75em}
  {}

\definecolor{mypurplemain}{RGB}{94,34,119}

\AddToShipoutPictureBG{%
\begin{tikzpicture}[remember picture,overlay]
\draw[mypurplemain, line width=2.5pt]
($(current page.north west)+(1cm,-1cm)$)
rectangle
($(current page.south east)+(-1cm,1cm)$);
\end{tikzpicture}%
}

\makeindex[intoc, title=Index]
\indexsetup{othercode={\lhead{\emph{Index}}}}

\ifx \nextra \undefined
\usepackage[pdftex,
hidelinks,
pdfauthor={Joseph Eddon},
pdfsubject={Durham Maths Notes: Year \npart\ - \ncourse},
pdftitle={Year \npart\ - \ncourse},
pdfkeywords={Durham Mathematics Maths Math \npart\ \nterm\ \nyear\ \ncourse}]{hyperref}
\title{Year \npart\ --- \ncourse}
\else
\usepackage[pdftex,
hidelinks,
pdfauthor={Joseph Eddon},
pdfsubject={Durham Maths Notes: Year \npart\ - \ncourse\ (\nextra)},
pdftitle={Year \npart\ - \ncourse\ (\nextra)},
pdfkeywords={Durham Mathematics Maths Math \npart\ \nterm\ \nyear\ \ncourse\ \nextra}]{hyperref}

\title{Year \npart\ --- \ncourse \\ {\Large \nextra}}
\renewcommand\printindex{}
\fi

\pgfplotsset{compat=1.12}

\pagestyle{fancyplain}

\setlength{\headsep}{20pt}

\ifx \ncoursehead \undefined
\def\ncoursehead{\ncourse}
\fi

\lhead{\emph{\nouppercase{\leftmark}}}
\ifx \nextra \undefined
\rhead{
\ifnum\thepage=1
\else
\npart\ \ncoursehead
\fi}
\else
\rhead{
\ifnum\thepage=1
\else
\npart\ \ncoursehead \ (\nextra)
\fi}
\fi
\usetikzlibrary{arrows.meta}
\usetikzlibrary{decorations.markings}
\usetikzlibrary{decorations.pathmorphing}
\usetikzlibrary{positioning}
\usetikzlibrary{fadings}
\usetikzlibrary{intersections}
\usetikzlibrary{cd}

\newcommand*{\Cdot}{{\raisebox{-0.25ex}{\scalebox{1.5}{$\cdot$}}}}
\newcommand {\pd}[2][ ]{
\ifx #1 { }
\frac{\partial}{\partial #2}
\else
\frac{\partial^{#1}}{\partial #2^{#1}}
\fi
}
\ifx \nhtml \undefined
\else
\renewcommand\printindex{}
\DisableLigatures[f]{family = *}
\let\Contentsline\contentsline
\renewcommand\contentsline[3]{\Contentsline{#1}{#2}{}}
\renewcommand{\@dotsep}{10000}
\newlength\currentparindent
\setlength\currentparindent\parindent

\newcommand\@minipagerestore{\setlength{\parindent}{\currentparindent}}
\usepackage[active,tightpage,pdftex]{preview}
\renewcommand{\PreviewBorder}{0.1cm}

\newenvironment{stretchpage}%
{\begin{preview}\begin{minipage}{\hsize}}%
{\end{minipage}\end{preview}}
\AtBeginDocument{\begin{stretchpage}}
\AtEndDocument{\end{stretchpage}}

\newcommand{\@@newpage}{\end{stretchpage}\begin{stretchpage}}

\let\@real@section\section
\renewcommand{\section}{\@@newpage\@real@section}
\let\@real@subsection\subsection
\renewcommand{\subsection}{\@ifstar{\@real@subsection*}{\@@newpage\@real@subsection}}
\fi
\ifx \ntrim \undefined
\else
\usepackage{geometry}
\geometry{
papersize={379pt, 699pt},
textwidth=345pt,
textheight=596pt,
left=17pt,
top=54pt,
right=17pt
}
\fi

\pdfinclusioncopyfonts=1

\setlength{\parindent}{0pt}
\setlength{\parskip}{0.5em}

% ============================================================
% Durham Purple Colour Palette
% ============================================================

% ------------------------------------------------------------
% Core Durham colours
% ------------------------------------------------------------

% Deep institutional purple
\definecolor{durhampurple}{RGB}{82,45,110}

% Rich violet
\definecolor{durhamviolet}{RGB}{108,68,150}

% Blue-violet
\definecolor{durhambluepurple}{RGB}{92,82,155}

% Soft lilac
\definecolor{durhamlilac}{RGB}{145,115,185}

% Muted mauve
\definecolor{durhammauve}{RGB}{160,95,145}

% Dark plum
\definecolor{durhamplum}{RGB}{105,55,100}


% Proof idea — soft green
\definecolor{proofideabg}{RGB}{230,242,232}
\definecolor{proofideaborder}{RGB}{55,115,75}

% Formal proof — neutral background
\definecolor{formalproofbg}{RGB}{245,245,248}
\definecolor{formalproofborder}{RGB}{100,100,115}

% ------------------------------------------------------------
% Box background colours
% ------------------------------------------------------------

% Definition — cool lavender
\definecolor{defbg}{RGB}{232,224,244}

% Theorem — blue
\definecolor{thmbg}{RGB}{211,222,244}

% Lemma / proposition — blue-lilac
\definecolor{lembg}{RGB}{231,229,246}

% Method / trick — soft mauve
\definecolor{methbg}{RGB}{244,226,238}

% Warning / trap — dusty plum-pink
\definecolor{warnbg}{RGB}{240,222,235}

% General supporting material
\definecolor{neutralbg}{RGB}{240,237,246}

\definecolor{thmborder}{RGB}{48,68,125}

% Soft cream question background
\definecolor{qbg}{RGB}{255, 240, 253}

% Question border
\definecolor{qborder}{RGB}{94,34,119}

% Black box — intentionally opaque / unexplained result
\definecolor{blackboxbg}{RGB}{35,35,38}

\definecolor{blackboxborder}{RGB}{80,80,85}

\definecolor{blackboxtext}{RGB}{245,245,245}

\definecolor{blackboxheading}{RGB}{210,210,215}

\definecolor{mistakebg}{RGB}{250,225,225}

\definecolor{mistakeborder}{RGB}{190,45,45}

\definecolor{mistaketext}{RGB}{55,25,25}

\definecolor{mistakeheading}{RGB}{150,25,25}


% ------------------------------------------------------------
% Box styles
% ------------------------------------------------------------

\newcommand{\boxfont}{\fontfamily{cmr}\selectfont}


% ============================================================
% Definitions
% ============================================================

\mdfdefinestyle{defstyle}{
    backgroundcolor=defbg,
    linecolor=durhamviolet,
    linewidth=1.0pt,
    roundcorner=4pt,
    skipabove=8pt,
    skipbelow=8pt,
    innertopmargin=7pt,
    innerbottommargin=7pt,
    innerleftmargin=8pt,
    innerrightmargin=8pt,
}


% ============================================================
% Theorems / methods / corollaries
% ============================================================

\mdfdefinestyle{thmstyle}{
    backgroundcolor=thmbg,
    linecolor=durhampurple,
    linewidth=1.0pt,
    roundcorner=4pt,
    skipabove=8pt,
    skipbelow=8pt,
    innertopmargin=7pt,
    innerbottommargin=7pt,
    innerleftmargin=8pt,
    innerrightmargin=8pt,
}


% ============================================================
% Lemmas / propositions / results
% ============================================================

\mdfdefinestyle{lemstyle}{
    backgroundcolor=lembg,
    linecolor=durhamlilac,
    linewidth=1.0pt,
    roundcorner=4pt,
    skipabove=8pt,
    skipbelow=8pt,
    innertopmargin=7pt,
    innerbottommargin=7pt,
    innerleftmargin=8pt,
    innerrightmargin=8pt,
}


% ============================================================
% Warnings / important observations
% ============================================================

\mdfdefinestyle{trapstyle}{
    backgroundcolor=warnbg,
    linecolor=durhamplum,
    linewidth=1.0pt,
    roundcorner=4pt,
    skipabove=8pt,
    skipbelow=8pt,
    innertopmargin=7pt,
    innerbottommargin=7pt,
    innerleftmargin=8pt,
    innerrightmargin=8pt,
}


% ============================================================
% Procedures / methods / tricks
% ============================================================

\mdfdefinestyle{doboxstyle}{
    backgroundcolor=methbg,
    linecolor=durhammauve,
    linewidth=1.0pt,
    roundcorner=4pt,
    skipabove=8pt,
    skipbelow=8pt,
    innertopmargin=7pt,
    innerbottommargin=7pt,
    innerleftmargin=8pt,
    innerrightmargin=8pt,
}

% ============================================================
% Questions
% ============================================================

\mdfdefinestyle{qstyle}{
    backgroundcolor=qbg,
    linecolor=qborder,
    linewidth=1.2pt,
    roundcorner=4pt,
    skipabove=10pt,
    skipbelow=10pt,
    innertopmargin=8pt,
    innerbottommargin=8pt,
    innerleftmargin=10pt,
    innerrightmargin=10pt,
}

% ============================================================
% Black Box
% ============================================================

\mdfdefinestyle{blackboxstyle}{
    backgroundcolor=blackboxbg,
    linecolor=blackboxborder,
    linewidth=1.0pt,
    roundcorner=4pt,
    skipabove=10pt,
    skipbelow=10pt,
    innertopmargin=9pt,
    innerbottommargin=9pt,
    innerleftmargin=10pt,
    innerrightmargin=10pt,
}

% ============================================================
% Mistakes
% ============================================================

\mdfdefinestyle{mistakestyle}{
    backgroundcolor=mistakebg,
    linecolor=mistakeborder,
    linewidth=1.0pt,
    roundcorner=4pt,
    skipabove=8pt,
    skipbelow=8pt,
    innertopmargin=7pt,
    innerbottommargin=7pt,
    innerleftmargin=8pt,
    innerrightmargin=8pt,
}


% ============================================================
% Proof Idea and Formal Proof
% ============================================================

\mdfdefinestyle{proofideastyle}{
    backgroundcolor=proofideabg,
    linecolor=proofideaborder,
    linewidth=1.0pt,
    roundcorner=4pt,
    skipabove=8pt,
    skipbelow=8pt,
    innertopmargin=7pt,
    innerbottommargin=7pt,
    innerleftmargin=8pt,
    innerrightmargin=8pt,
    nobreak=false,
    splittopskip=10pt,
    splitbottomskip=10pt,
}

\mdfdefinestyle{formalproofstyle}{
    backgroundcolor=formalproofbg,
    linecolor=formalproofborder,
    linewidth=1.0pt,
    roundcorner=4pt,
    skipabove=8pt,
    skipbelow=8pt,
    innertopmargin=7pt,
    innerbottommargin=7pt,
    innerleftmargin=8pt,
    innerrightmargin=8pt,
    nobreak=false,
    splittopskip=10pt,
    splitbottomskip=10pt,
}


% Callout environments

% ============================================================
% Theorem environments
% ============================================================

% ------------------------------------------------------------
% Numbered theorem counter
% ------------------------------------------------------------

\newtheorem{genericthm}{Result}[section]


% ============================================================
% Theorem heading styles
% ============================================================

\newtheoremstyle{defnstyle}
{6pt}
{6pt}
{\normalfont}
{0pt}
{\bfseries\color{durhamviolet}}
{.}
{0.5em}
{}

\newtheoremstyle{thmhead}
{6pt}
{6pt}
{\normalfont}
{0pt}
{\bfseries\color{thmborder}}
{.}
{0.5em}
{}

\newtheoremstyle{lemhead}
{6pt}
{6pt}
{\normalfont}
{0pt}
{\bfseries\color{durhamlilac!80!black}}
{.}
{0.5em}
{}

\newtheoremstyle{plainhead}
{6pt}
{6pt}
{\normalfont}
{0pt}
{\bfseries\color{durhamplum}}
{.}
{0.5em}
{}
% ------------------------------------------------------------
% Internal theorem environments
% ------------------------------------------------------------

\theoremstyle{defnstyle}
\newtheorem{definition-inner}[genericthm]{Definition}

\theoremstyle{thmhead}
\newtheorem{theorem-inner}[genericthm]{Theorem}
\newtheorem{corollary-inner}[genericthm]{Corollary}
\newtheorem{method-inner}[genericthm]{Method}
\newtheorem{trick-inner}[genericthm]{Trick}

\theoremstyle{lemhead}
\newtheorem{lemma-inner}[genericthm]{Lemma}
\newtheorem{proposition-inner}[genericthm]{Proposition}
\newtheorem{result-inner}[genericthm]{Result}

% ============================================================
% Question command
% ============================================================

\newcommand{\qs}[2]{%
    \begin{mdframed}[style=qstyle]
        \textbf{\color{qborder}#1}\par
        \vspace{4pt}
        #2
    \end{mdframed}
}

% ------------------------------------------------------------
% Public boxed environments
% ------------------------------------------------------------

\newenvironment{definition}[1][]%
{%
\begin{mdframed}[style=defstyle]
\begin{definition-inner}[#1]%
}%
{%
\end{definition-inner}
\end{mdframed}%
}

\newenvironment{theorem}[1][]%
{%
\begin{mdframed}[style=thmstyle]
\begin{theorem-inner}[#1]%
}%
{%
\end{theorem-inner}
\end{mdframed}%
}

\newenvironment{method}[1][]%
{%
\begin{mdframed}[style=thmstyle]
\begin{method-inner}[#1]%
}%
{%
\end{method-inner}
\end{mdframed}%
}

\newenvironment{trick}[1][]%
{%
\begin{mdframed}[style=thmstyle]
\begin{trick-inner}[#1]%
}%
{%
\end{trick-inner}
\end{mdframed}%
}

\newenvironment{corollary}[1][]%
{%
\begin{mdframed}[style=thmstyle]
\begin{corollary-inner}[#1]%
}%
{%
\end{corollary-inner}
\end{mdframed}%
}

\newenvironment{lemma}[1][]%
{%
\begin{mdframed}[style=lemstyle]
\begin{lemma-inner}[#1]%
}%
{%
\end{lemma-inner}
\end{mdframed}%
}

\newenvironment{proposition}[1][]%
{%
\begin{mdframed}[style=lemstyle]
\begin{proposition-inner}[#1]%
}%
{%
\end{proposition-inner}
\end{mdframed}%
}

\newenvironment{result}[1][]%
{%
\begin{mdframed}[style=lemstyle]
\begin{result-inner}[#1]%
}%
{%
\end{result-inner}
\end{mdframed}%
}

\newenvironment{blackbox}[1][]%
{%
    \begin{mdframed}[style=blackboxstyle]
    \color{blackboxtext}
    \ifx\relax#1\relax
        \textbf{\color{blackboxheading}Black Box}\par
    \else
        \textbf{\color{blackboxheading}#1}\par
    \fi
    \vspace{4pt}
}%
{%
    \end{mdframed}
}

\newenvironment{mistake}[1][]%
{%
\begin{mdframed}[style=mistakestyle]
\textbf{\color{mistakeborder}%
\ifx\relax#1\relax
Mistake%
\else
#1%
\fi
}\par
\vspace{4pt}
}%
{%
\end{mdframed}%
}

% ============================================================
% Proof Idea + Formal Proof Environment
% ============================================================

\newenvironment{proofidea}
{%
    \par\addvspace{6pt}
    \begin{mdframed}[style=proofideastyle]
    \textbf{\color{proofideaborder}Main Idea / Strategy}
    \par\vspace{4pt}
}
{%
    \end{mdframed}
}

\newenvironment{formalproof}
{%
    \begin{mdframed}[style=formalproofstyle]
    \textbf{\color{formalproofborder}Formal Proof}
    \par\vspace{4pt}
    \pushQED{\qed}
}
{%
    \popQED
    \end{mdframed}
}

% ------------------------------------------------------------
% Unboxed numbered environments
% ------------------------------------------------------------

\theoremstyle{plainhead}

\newtheorem{example}[genericthm]{Example}
\newtheorem{remark}[genericthm]{Remark}
\newtheorem{notation}[genericthm]{Notation}



% ------------------------------------------------------------
% Compatibility environments
% ------------------------------------------------------------

% n-prefixed versions

\newenvironment{ndefi}[1][]%
{\begin{definition}[#1]}%
{\end{definition}}

\newenvironment{nthm}[1][]%
{\begin{theorem}[#1]}%
{\end{theorem}}

\newenvironment{nlem}[1][]%
{\begin{lemma}[#1]}%
{\end{lemma}}

\newenvironment{nprop}[1][]%
{\begin{proposition}[#1]}%
{\end{proposition}}

\newenvironment{ncor}[1][]%
{\begin{corollary}[#1]}%
{\end{corollary}}

\newenvironment{nmeth}[1][]%
{\begin{method}[#1]}%
{\end{method}}

\newenvironment{ntfn}[1][]%
{\begin{trick}[#1]}%
{\end{trick}}

\newenvironment{nres}[1][]%
{\begin{result}[#1]}%
{\end{result}}

% Examples / remarks

\newenvironment{nex}[1][]%
{\begin{example}[#1]}%
{\end{example}}

\newenvironment{nrem}[1][]%
{\begin{remark}[#1]}%
{\end{remark}}


% Short versions

\newenvironment{defi}[1][]%
{\begin{definition}[#1]}%
{\end{definition}}

\newenvironment{thm}[1][]%
{\begin{theorem}[#1]}%
{\end{theorem}}

\newenvironment{prop}[1][]%
{\begin{proposition}[#1]}%
{\end{proposition}}

\newenvironment{cor}[1][]%
{\begin{corollary}[#1]}%
{\end{corollary}}

\newenvironment{eg}[1][]%
{\begin{example}[#1]}%
{\end{example}}

\newenvironment{rem}[1][]%
{\begin{remark}[#1]}%
{\end{remark}}

\newenvironment{res}[1][]%
{\begin{result}[#1]}%
{\end{result}}

\newenvironment{meth}[1][]%
{\begin{method}[#1]}%
{\end{method}}

\newenvironment{tfn}[1][]%
{\begin{trick}[#1]}%
{\end{trick}}


\renewcommand{\labelitemi}{--}
\renewcommand{\labelitemii}{$\circ$}
\renewcommand{\labelenumi}{(\roman{*})}

%\let\stdsection\section
%\renewcommand\section{\newpage\stdsection}



\newcommand\qedsym{\hfill\ensuremath{\square}}
% Strike through
\def\st{\bgroup \ULdepth=-.55ex \ULset}


%%%%%%%%%%%%%%%%%%%%%%%%%
%%%%% Maths Symbols %%%%%
%%%%%%%%%%%%%%%%%%%%%%%%%

% Matrix groups
\newcommand{\GL}{\mathrm{GL}}
\newcommand{\Or}{\mathrm{O}}
\newcommand{\PGL}{\mathrm{PGL}}
\newcommand{\PSL}{\mathrm{PSL}}
\newcommand{\PSO}{\mathrm{PSO}}
\newcommand{\PSU}{\mathrm{PSU}}
\newcommand{\SL}{\mathrm{SL}}
\newcommand{\SO}{\mathrm{SO}}
\newcommand{\Spin}{\mathrm{Spin}}
\newcommand{\Sp}{\mathrm{Sp}}
\newcommand{\SU}{\mathrm{SU}}
\newcommand{\U}{\mathrm{U}}
\newcommand{\Mat}{\mathrm{Mat}}

% Matrix algebras
\newcommand{\gl}{\mathfrak{gl}}
\newcommand{\ort}{\mathfrak{o}}
\newcommand{\so}{\mathfrak{so}}
\newcommand{\su}{\mathfrak{su}}
\newcommand{\uu}{\mathfrak{u}}
\renewcommand{\sl}{\mathfrak{sl}}

% Special sets
\newcommand{\C}{\mathbb{C}}
\newcommand{\CP}{\mathbb{CP}}
\newcommand{\GG}{\mathbb{G}}
\newcommand{\N}{\mathbb{N}}
\newcommand{\Q}{\mathbb{Q}}
\newcommand{\R}{\mathbb{R}}
\newcommand{\RP}{\mathbb{RP}}
\newcommand{\T}{\mathbb{T}}
\newcommand{\Z}{\mathbb{Z}}
\renewcommand{\H}{\mathbb{H}}

% Brackets
\newcommand{\abs}[1]{\left\lvert #1\right\rvert}
\newcommand{\bket}[1]{\left\lvert #1\right\rangle}
\newcommand{\brak}[1]{\left\langle #1 \right\rvert}
\newcommand{\braket}[2]{\left\langle #1\middle\vert #2 \right\rangle}
\newcommand{\bra}{\langle}
\newcommand{\ket}{\rangle}
\newcommand{\norm}[1]{\left\lVert #1\right\rVert}
\newcommand{\normalorder}[1]{\mathop{:}\nolimits\!#1\!\mathop{:}\nolimits}
\newcommand{\tv}[1]{|#1|}
\renewcommand{\vec}[1]{\boldsymbol{\mathbf{#1}}}

% not-math
\newcommand{\bolds}[1]{{\bfseries #1}}
\newcommand{\cat}[1]{\mathsf{#1}}
\newcommand{\ph}{\,\cdot\,}
\newcommand{\term}[1]{\emph{#1}\index{#1}}
\newcommand{\phantomeq}{\hphantom{{}={}}}
% Probability
\DeclareMathOperator{\Bernoulli}{Bernoulli}
\DeclareMathOperator{\betaD}{beta}
\DeclareMathOperator{\bias}{bias}
\DeclareMathOperator{\binomial}{binomial}
\DeclareMathOperator{\corr}{corr}
\DeclareMathOperator{\cov}{cov}
\DeclareMathOperator{\gammaD}{gamma}
\DeclareMathOperator{\mse}{mse}
\DeclareMathOperator{\multinomial}{multinomial}
\DeclareMathOperator{\Poisson}{Poisson}
\DeclareMathOperator{\var}{var}
\newcommand{\E}{\mathbb{E}}
\newcommand{\Prob}{\mathbb{P}}

% Algebra
\DeclareMathOperator{\adj}{adj}
\DeclareMathOperator{\Ann}{Ann}
\DeclareMathOperator{\Aut}{Aut}
\DeclareMathOperator{\Char}{char}
\DeclareMathOperator{\disc}{disc}
\DeclareMathOperator{\dom}{dom}
\DeclareMathOperator{\fix}{fix}
\DeclareMathOperator{\Hom}{Hom}
\DeclareMathOperator{\id}{id}
\DeclareMathOperator{\image}{image}
\DeclareMathOperator{\im}{im}
\DeclareMathOperator{\tr}{tr}
\DeclareMathOperator{\Tr}{Tr}
\newcommand{\Bilin}{\mathrm{Bilin}}
\newcommand{\Frob}{\mathrm{Frob}}

% Others
\newcommand\ad{\mathrm{ad}}
\newcommand\Art{\mathrm{Art}}
\newcommand{\B}{\mathcal{B}}
\newcommand{\cU}{\mathcal{U}}
\newcommand{\Der}{\mathrm{Der}}
\newcommand{\D}{\mathrm{D}}
\newcommand{\dR}{\mathrm{dR}}
\newcommand{\exterior}{\mathchoice{{\textstyle\bigwedge}}{{\bigwedge}}{{\textstyle\wedge}}{{\scriptstyle\wedge}}}
\newcommand{\F}{\mathbb{F}}
\newcommand{\G}{\mathcal{G}}
\newcommand{\Gr}{\mathrm{Gr}}
\newcommand{\haut}{\mathrm{ht}}
\newcommand{\Hol}{\mathrm{Hol}}
\newcommand{\hol}{\mathfrak{hol}}
\newcommand{\Id}{\mathrm{Id}}
\newcommand{\lie}[1]{\mathfrak{#1}}
\newcommand{\op}{\mathrm{op}}
\newcommand{\Oc}{\mathcal{O}}
\newcommand{\pr}{\mathrm{pr}}
\newcommand{\Ps}{\mathcal{P}}
\newcommand{\pt}{\mathrm{pt}}
\newcommand{\qeq}{\mathrel{``{=}"}}
\newcommand{\Rs}{\mathcal{R}}
\newcommand{\Vect}{\mathrm{Vect}}
\newcommand{\wsto}{\stackrel{\mathrm{w}^*}{\to}}
\newcommand{\wt}{\mathrm{wt}}
\newcommand{\wto}{\stackrel{\mathrm{w}}{\to}}
\renewcommand{\d}{\mathrm{d}}
\renewcommand{\P}{\mathbb{P}}
\newcommand{\nthmno}[1]{%
\par\addvspace{\baselineskip} % Adds one full line of space
\noindent
This is \textbf{Theorem~#1} from the official lecture notes.
\par\addvspace{\baselineskip}
}

\newcommand{\ndefno}[1]{%
\par\addvspace{\baselineskip} % Adds one full line of space
\noindent
This is \textbf{Definition~#1} from the official lecture notes.
\par\addvspace{\baselineskip}
}

\newcommand{\npropno}[1]{%
\par\addvspace{\baselineskip} % Adds one full line of space
\noindent
This is \textbf{Proposition~#1} from the official lecture notes.
\par\addvspace{\baselineskip}
}

\newcommand{\ncorno}[1]{%
\par\addvspace{\baselineskip} % Adds one full line of space
\noindent
This is \textbf{Corollary~#1} from the official lecture notes.
\par\addvspace{\baselineskip}
}

\newcommand{\nlemno}[1]{%
\par\addvspace{\baselineskip} % Adds one full line of space
\noindent
This is \textbf{Lemma~#1} from the official lecture notes.
\par\addvspace{\baselineskip}
}


%\renewcommand{\F}{\mathcal{F}}


\let\Im\relax
\let\Re\relax

\DeclareMathOperator{\area}{area}
\DeclareMathOperator{\card}{card}
\DeclareMathOperator{\ccl}{ccl}
\DeclareMathOperator{\ch}{ch}
\DeclareMathOperator{\cl}{cl}
\DeclareMathOperator{\cls}{\overline{\mathrm{span}}}
\DeclareMathOperator{\coker}{coker}
\DeclareMathOperator{\conv}{conv}
\DeclareMathOperator{\cosec}{cosec}
\DeclareMathOperator{\cosech}{cosech}
\DeclareMathOperator{\covol}{covol}
\DeclareMathOperator{\diag}{diag}
\DeclareMathOperator{\diam}{diam}
\DeclareMathOperator{\Diff}{Diff}
\DeclareMathOperator{\End}{End}
\DeclareMathOperator{\energy}{energy}
\DeclareMathOperator{\erfc}{erfc}
\DeclareMathOperator{\erf}{erf}
\DeclareMathOperator*{\esssup}{ess\,sup}
\DeclareMathOperator{\ev}{ev}
\DeclareMathOperator{\Ext}{Ext}
\DeclareMathOperator{\fst}{fst}
\DeclareMathOperator{\Fit}{Fit}
\DeclareMathOperator{\Frac}{Frac}
\DeclareMathOperator{\Gal}{Gal}
\DeclareMathOperator{\gr}{gr}
\DeclareMathOperator{\hcf}{hcf}
\DeclareMathOperator{\Im}{Im}
\DeclareMathOperator{\Ind}{Ind}
\DeclareMathOperator{\Int}{Int}
\DeclareMathOperator{\Isom}{Isom}
\DeclareMathOperator{\lcm}{lcm}
\DeclareMathOperator{\length}{length}
\DeclareMathOperator{\Lie}{Lie}
\DeclareMathOperator{\like}{like}
\DeclareMathOperator{\Lk}{Lk}
\DeclareMathOperator{\Maps}{Maps}
\DeclareMathOperator{\orb}{orb}
\DeclareMathOperator{\ord}{ord}
\DeclareMathOperator{\otp}{otp}
\DeclareMathOperator{\poly}{poly}
\DeclareMathOperator{\rank}{rank}
\DeclareMathOperator{\rel}{rel}
\DeclareMathOperator{\Rad}{Rad}
\DeclareMathOperator{\Re}{Re}
\DeclareMathOperator{\Res}{Res}
\DeclareMathOperator{\Ric}{Ric}
\DeclareMathOperator{\rk}{rk}
\DeclareMathOperator{\Rees}{Rees}
\DeclareMathOperator{\Root}{Root}
\DeclareMathOperator{\sech}{sech}
\DeclareMathOperator{\sgn}{sgn}
\DeclareMathOperator{\snd}{snd}
\DeclareMathOperator{\Spec}{Spec}
\DeclareMathOperator{\spn}{span}
\DeclareMathOperator{\stab}{stab}
\DeclareMathOperator{\St}{St}
\DeclareMathOperator{\supp}{supp}
\DeclareMathOperator{\Syl}{Syl}
\DeclareMathOperator{\Sym}{Sym}
\DeclareMathOperator{\vol}{vol}

\pgfarrowsdeclarecombine{twolatex'}{twolatex'}{latex'}{latex'}{latex'}{latex'}
\tikzset{->/.style = {decoration={markings,
mark=at position 1 with {\arrow[scale=2]{latex'}}},
postaction={decorate}}}
\tikzset{<-/.style = {decoration={markings,
mark=at position 0 with {\arrowreversed[scale=2]{latex'}}},
postaction={decorate}}}
\tikzset{<->/.style = {decoration={markings,
mark=at position 0 with {\arrowreversed[scale=2]{latex'}},
mark=at position 1 with {\arrow[scale=2]{latex'}}},
postaction={decorate}}}
\tikzset{->-/.style = {decoration={markings,
mark=at position #1 with {\arrow[scale=2]{latex'}}},
postaction={decorate}}}
\tikzset{-<-/.style = {decoration={markings,
mark=at position #1 with {\arrowreversed[scale=2]{latex'}}},
postaction={decorate}}}
\tikzset{->>/.style = {decoration={markings,
mark=at position 1 with {\arrow[scale=2]{latex'}}},
postaction={decorate}}}
\tikzset{<<-/.style = {decoration={markings,
mark=at position 0 with {\arrowreversed[scale=2]{twolatex'}}},
postaction={decorate}}}
\tikzset{<<->>/.style = {decoration={markings,
mark=at position 0 with {\arrowreversed[scale=2]{twolatex'}},
mark=at position 1 with {\arrow[scale=2]{twolatex'}}},
postaction={decorate}}}
\tikzset{->>-/.style = {decoration={markings,
mark=at position #1 with {\arrow[scale=2]{twolatex'}}},
postaction={decorate}}}
\tikzset{-<<-/.style = {decoration={markings,
mark=at position #1 with {\arrowreversed[scale=2]{twolatex'}}},
postaction={decorate}}}

\tikzset{circ/.style = {fill, circle, inner sep = 0, minimum size = 3}}
\tikzset{scirc/.style = {fill, circle, inner sep = 0, minimum size = 1.5}}
\tikzset{mstate/.style={circle, draw, blue, text=black, minimum width=0.7cm}}

\tikzset{eqpic/.style={baseline={([yshift=-.5ex]current bounding box.center)}}}
\tikzset{commutative diagrams/.cd,cdmap/.style={/tikz/column 1/.append style={anchor=base east},/tikz/column 2/.append style={anchor=base west},row sep=tiny}}

\definecolor{mblue}{rgb}{0.2, 0.3, 0.8}
\definecolor{morange}{rgb}{1, 0.5, 0}
\definecolor{mgreen}{rgb}{0.1, 0.4, 0.2}
\definecolor{mred}{rgb}{0.5, 0, 0}

\def\drawcirculararc(#1,#2)(#3,#4)(#5,#6){%
\pgfmathsetmacro\cA{(#1*#1+#2*#2-#3*#3-#4*#4)/2}%
\pgfmathsetmacro\cB{(#1*#1+#2*#2-#5*#5-#6*#6)/2}%
\pgfmathsetmacro\cy{(\cB*(#1-#3)-\cA*(#1-#5))/%
((#2-#6)*(#1-#3)-(#2-#4)*(#1-#5))}%
\pgfmathsetmacro\cx{(\cA-\cy*(#2-#4))/(#1-#3)}%
\pgfmathsetmacro\cr{sqrt((#1-\cx)*(#1-\cx)+(#2-\cy)*(#2-\cy))}%
\pgfmathsetmacro\cA{atan2(#2-\cy,#1-\cx)}%
\pgfmathsetmacro\cB{atan2(#6-\cy,#5-\cx)}%
\pgfmathparse{\cB<\cA}%
\ifnum\pgfmathresult=1
\pgfmathsetmacro\cB{\cB+360}%
\fi
\draw (#1,#2) arc (\cA:\cB:\cr);%
}
\newcommand\getCoord[3]{\newdimen{#1}\newdimen{#2}\pgfextractx{#1}{\pgfpointanchor{#3}{center}}\pgfextracty{#2}{\pgfpointanchor{#3}{center}}}

\newcommand\qedshift{\vspace{-17pt}}
\newcommand\fakeqed{\pushQED{\qed}\qedhere}

\def\Xint#1{\mathchoice
{\XXint\displaystyle\textstyle{#1}}%
{\XXint\textstyle\scriptstyle{#1}}%
{\XXint\scriptstyle\scriptscriptstyle{#1}}%
{\XXint\scriptscriptstyle\scriptscriptstyle{#1}}%
\!\int}
\def\XXint#1#2#3{{\setbox0=\hbox{$#1{#2#3}{\int}$}
\vcenter{\hbox{$#2#3$}}\kern-.5\wd0}}
\def\ddashint{\Xint=}
\def\dashint{\Xint-}

\newcommand\separator{{\centering\rule{2cm}{0.2pt}\vspace{2pt}\par}}

\newenvironment{own}{\color{gray!70!black}}{}

\newcommand\makecenter[1]{\raisebox{-0.5\height}{#1}}

\mathchardef\mdash="2D

\newenvironment{significant}{\begin{center}\begin{minipage}{0.9\textwidth}\centering\em}{\end{minipage}\end{center}}
\DeclareRobustCommand{\rvdots}{%
\vbox{
\baselineskip4\p@\lineskiplimit\z@
\kern-\p@
\hbox{.}\hbox{.}\hbox{.}
}}
\DeclareRobustCommand\tph[3]{{\texorpdfstring{#1}{#2}}}
\makeatother
